\documentclass[conference]{IEEEtran}

\usepackage{cite}
\usepackage{amsmath,amssymb,amsfonts}
\usepackage{graphicx}
\usepackage{textcomp}
\usepackage{booktabs}
\usepackage{array}
\usepackage{siunitx}
\usepackage{hyperref}
\usepackage{listings}
\usepackage{xcolor}

\lstset{
  language=Matlab,
  basicstyle=\ttfamily\footnotesize,
  keywordstyle=\color{blue},
  commentstyle=\color{green!50!black},
  stringstyle=\color{red!60!black},
  numbers=left,
  numberstyle=\tiny\color{gray},
  frame=single,
  breaklines=true,
  captionpos=b
}

\def\BibTeX{{\rm B\kern-.05em{\sc i\kern-.025em b}\kern-.08em T\kern-.1667em\lower.7ex\hbox{E}\kern-.125emX}}

\begin{document}

\title{Modular Behavioral Modeling and Signal Integrity Analysis of a UCIe 3.0 Transmitter--Receiver Link at 32/64~GT/s}

\author{\IEEEauthorblockN{Sai Prakash}
\IEEEauthorblockA{2nd Sem, MTech in SOCD\\
Indian Institute of Technology Palakkad\\
Kerala, India\\
152502010@smail.iitpkd.ac.in}
\and
\IEEEauthorblockN{Pushpa Rani M P}
\IEEEauthorblockA{2nd Sem, MTech in SOCD\\
Indian Institute of Technology Palakkad\\
Kerala, India\\
152502015@smail.iitpkd.ac.in}
}

\maketitle

\begin{abstract}
This paper presents a comprehensive, modular behavioral model for a Universal Chiplet Interconnect Express (UCIe) 3.0 high-speed serial link operating at 32 and 64~GT/s. The complete transmitter-to-receiver signal path is decomposed into twelve independently executable functional blocks implemented in MATLAB, spanning PRBS-13 pattern generation, NRZ mapping, 2-tap feed-forward equalization (FFE), analog driver shaping with Gaussian rise/fall modeling, RC channel loss, receiver termination, additive noise injection, clock-jitter-aware sampling, threshold slicing, direct BER measurement, Q-factor estimation via voltage histograms, and eye diagram accumulation. Each block is encapsulated as a parameterized function with a standardized interface, enabling rapid reconfiguration for diverse signal integrity scenarios. Five industry-representative test cases are developed: (i)~advanced silicon interposer ultra-short reach, (ii)~standard organic package medium reach, (iii)~next-generation 64~GT/s bandwidth-limited operation, (iv)~severe crosstalk and power supply noise stress, and (v)~PLL clock jitter failure analysis. The modular architecture enables independent block-level validation and automated plot generation for each processing stage. Simulation results demonstrate that the framework accurately captures the interplay between channel parasitics, equalization strength, voltage noise, and timing jitter, providing a flexible platform for UCIe link design exploration and signal integrity education.
\end{abstract}

\begin{IEEEkeywords}
UCIe 3.0, behavioral modeling, signal integrity, feed-forward equalization, eye diagram, bit error rate, high-speed serial link, chiplet interconnect
\end{IEEEkeywords}

%% ====================================================================
\section{Introduction}
%% ====================================================================

The Universal Chiplet Interconnect Express (UCIe) standard has emerged as the dominant die-to-die interconnect specification for heterogeneous chiplet integration, targeting data rates of 32~GT/s in the current generation and 64~GT/s in future revisions~\cite{ucie_spec}. As multi-die packaging architectures---including 2.5D silicon interposers, fan-out wafer-level packaging (FOWLP), and organic substrates---become mainstream in advanced SoC designs, accurate behavioral modeling of the complete link becomes essential for early-stage signal integrity (SI) assessment~\cite{razavi2012}.

Traditional full-wave electromagnetic simulations, while highly accurate, are computationally prohibitive for system-level link analysis involving thousands of unit intervals (UI). Behavioral models offer a practical alternative by abstracting the physical phenomena into parameterized transfer functions that execute orders of magnitude faster while retaining sufficient accuracy for design-space exploration~\cite{fan2007signal}.

This work presents a fully modular MATLAB-based behavioral model of a UCIe 3.0 link. The key contributions are:

\begin{enumerate}
    \item \textbf{12-Block Modular Architecture:} The link is decomposed into twelve independently executable blocks spanning the transmitter, channel, receiver, and measurement stages. Each block follows a standardized function interface \texttt{[output, plotData] = block\_xxx(input, cfg)}, enabling plug-and-play substitution.
    
    \item \textbf{Physical Parameterization:} All RC parasitic values, FFE tap coefficients, noise levels, and jitter budgets are drawn from realistic UCIe specifications and can be swept via a centralized configuration structure.
    
    \item \textbf{Industry-Representative Test Cases:} Five test cases model real-world packaging scenarios---from premium silicon interposers to degraded PLL clock recovery---demonstrating vertical eye closure, horizontal eye closure, and the role of equalization.
    
    \item \textbf{Automated Visualization:} Each block produces per-stage diagnostic plots that are automatically saved as PNG images, enabling rapid visual inspection of signal degradation through the link.
\end{enumerate}

%% ====================================================================
\section{System Architecture}
%% ====================================================================

The complete link model follows the signal flow shown in Fig.~\ref{fig:arch}. The twelve blocks are organized into four sections.

\begin{figure}[h]
\centering
\fbox{\parbox{0.95\columnwidth}{\centering\small
\textbf{Section 1: Transmitter}\\
PRBS-13 $\rightarrow$ NRZ Mapper $\rightarrow$ TX FFE $\rightarrow$ TX Driver\\[4pt]
$\downarrow$\\[2pt]
\textbf{Section 2: Channel \& RX Load}\\
Channel RC $\rightarrow$ RX Termination\\[4pt]
$\downarrow$\\[2pt]
\textbf{Section 3: Receiver Sampler}\\
Add Noise $\rightarrow$ Sample \& Hold $\rightarrow$ Slicer\\[4pt]
$\downarrow$\\[2pt]
\textbf{Section 4: Measurement}\\
BER $\quad$ Histogram/Q $\quad$ Eye Diagram
}}
\caption{Complete UCIe link behavioral model architecture with twelve functional blocks organized into four processing sections.}
\label{fig:arch}
\end{figure}

\subsection{Design Philosophy}

The modular design follows three principles:

\begin{itemize}
    \item \textbf{Separation of Concerns:} Each block encapsulates exactly one physical phenomenon (e.g., channel loss, noise injection), making the model easy to extend and debug.
    \item \textbf{Standalone Executability:} Every block includes a \texttt{nargin==0} self-test that automatically generates upstream dependencies with default parameters, allowing any block to be run independently.
    \item \textbf{Centralized Configuration:} A single \texttt{cfg} structure propagates all parameters through the chain, enabling one-line test case definition.
\end{itemize}

%% ====================================================================
\section{Transmitter Modeling (Section 1)}
%% ====================================================================

\subsection{Block 1: PRBS-13 Generator}

The stimulus is generated using a 13-bit Linear Feedback Shift Register (LFSR) with feedback polynomial:
\begin{equation}
    P(x) = x^{13} + x^4 + x^3 + x + 1
\end{equation}
producing a maximal-length sequence of period $2^{13}-1 = 8191$ bits. The PRBS-13 pattern provides adequate spectral coverage for link characterization at data rates up to 64~GT/s while maintaining a manageable sequence length~\cite{jeruchim2000}.

\subsection{Block 2: NRZ Mapper}

Binary bits $b[n] \in \{0, 1\}$ are mapped to bipolar NRZ symbols:
\begin{equation}
    x[n] = 2 \cdot b[n] - 1 \in \{-1, +1\}
\end{equation}

\subsection{Block 3: TX Feed-Forward Equalizer (FFE)}

The UCIe specification defines a 2-tap FFE operating at the symbol rate. The equalized output is:
\begin{equation}
    y[n] = c_0 \cdot x[n] + c_1 \cdot x[n-1]
\end{equation}
where $c_0$ is the main cursor tap and $c_1$ is the post-cursor de-emphasis tap. The constraint $|c_0| + |c_1| = 1$ ensures unit energy normalization. Four standard presets are defined in Table~\ref{tab:presets}.

\begin{table}[h]
\caption{UCIe FFE Tap Presets}
\label{tab:presets}
\centering
\begin{tabular}{|c|c|c|l|}
\hline
\textbf{Preset} & $c_0$ & $c_1$ & \textbf{Description} \\
\hline
0 & 1.00 & 0.00 & No equalization \\
1 & 0.90 & $-$0.10 & Mild de-emphasis \\
2 & 0.85 & $-$0.15 & Moderate de-emphasis \\
3 & 0.75 & $-$0.25 & Strong de-emphasis \\
\hline
\end{tabular}
\end{table}

The FFE pre-distorts the transmitted waveform by attenuating the main cursor and boosting high-frequency transitions. This compensates for the low-pass filtering effect of the channel, effectively ``opening'' the eye at the receiver~\cite{proakis2008}.

\subsection{Block 4: TX Driver}

The driver block consolidates three analog operations:

\subsubsection{Zero-Order Hold (ZOH) Upsampling}
Each symbol is held for $N = 16$ samples, producing a waveform at sampling rate $f_s = N \times f_{\text{data}}$:
\begin{equation}
    f_s = 16 \times 32 \times 10^9 = 512 \text{ GHz}
\end{equation}

\subsubsection{Gaussian Rise/Fall Shaping}
Physical CMOS drivers cannot produce instantaneous voltage transitions. The finite slew rate is modeled by convolving the ideal waveform with a Gaussian kernel:
\begin{equation}
    g[k] = \frac{1}{\sqrt{2\pi}\sigma_s} \exp\!\left(-\frac{k^2}{2\sigma_s^2}\right)
\end{equation}
where $\sigma_s = t_{r,20\text{-}80} / 1.6832$ converts the 20\%--80\% rise time target to the Gaussian standard deviation. For $t_r = 6.25$~ps at 32~GT/s, $\sigma_s \approx 3.71$~ps.

\subsubsection{Output Impedance RC}
The TX pad introduces a series resistance $R_{\text{drv}} = 30~\Omega$ and parasitic shunt capacitance $C_{\text{drv}} = 0.1$~pF, modeled as a first-order IIR filter:
\begin{equation}
    \alpha = e^{-1/(f_s \cdot R_{\text{drv}} C_{\text{drv}})}, \quad H(z) = \frac{1-\alpha}{1 - \alpha z^{-1}}
\end{equation}

%% ====================================================================
\section{Channel and RX Load Modeling (Section 2)}
%% ====================================================================

\subsection{Block 5: Channel RC Model}

The die-to-die interconnect trace is modeled as a single-pole RC low-pass filter with configurable resistance $R_{\text{ch}}$ and capacitance $C_{\text{ch}}$. The 3~dB bandwidth is:
\begin{equation}
    f_{3\text{dB}} = \frac{1}{2\pi R_{\text{ch}} C_{\text{ch}}}
\end{equation}

The continuous-time transfer function $H(s) = \omega_0 / (s + \omega_0)$ is converted to discrete time via the bilinear transform implemented manually (without requiring the Signal Processing Toolbox):
\begin{equation}
    s = \frac{2f_s (z-1)}{z+1}, \quad H(z) = \frac{\omega_0(1 + z^{-1})}{(2f_s + \omega_0) + ({\omega_0 - 2f_s})z^{-1}}
\end{equation}

This parameterization allows different packaging technologies to be modeled by varying $R_{\text{ch}}$ and $C_{\text{ch}}$---from premium silicon interposers ($R=15~\Omega$, $C=0.05$~pF) to lossy organic substrates ($R=40~\Omega$, $C=0.15$~pF).

\subsection{Block 6: RX Termination Load}

The receiver pad presents a parallel termination of $R_{\text{rx}} = 50~\Omega$ and $C_{\text{rx}} = 0.2$~pF to ground. The $50~\Omega$ resistor absorbs the incident signal energy to suppress reflections at the far end of the transmission line, while $C_{\text{rx}}$ models the parasitic pad capacitance:
\begin{equation}
    \tau_{\text{rx}} = R_{\text{rx}} \cdot C_{\text{rx}} = 50 \times 0.2 \times 10^{-12} = 10 \text{ ps}
\end{equation}

%% ====================================================================
\section{Receiver Sampler Modeling (Section 3)}
%% ====================================================================

\subsection{Block 7: Additive Noise}

Real-world voltage noise from crosstalk, power supply ripple, and thermal sources is modeled as Additive White Gaussian Noise (AWGN):
\begin{equation}
    r_{\text{noisy}}[k] = r_{\text{clean}}[k] + n[k], \quad n[k] \sim \mathcal{N}(0, \sigma_V^2)
\end{equation}
where $\sigma_V$ ranges from 1~mV (clean environment) to 20~mV (severe crosstalk) depending on the test scenario.

\subsection{Block 8: Sample and Hold}

The receiver clock recovery circuit samples the waveform once per UI at the center of the eye. The ideal sampling instant for bit $i$ is:
\begin{equation}
    t_{\text{ideal}}[i] = \left(i - \tfrac{1}{2}\right) \cdot T_{\text{UI}} + \tfrac{T_{\text{UI}}}{2}
\end{equation}

When clock jitter is enabled, a random timing perturbation is added:
\begin{equation}
    t_{\text{actual}}[i] = t_{\text{ideal}}[i] + \Delta t_j, \quad \Delta t_j \sim \mathcal{N}(0, \sigma_{\text{RJ}}^2)
\end{equation}
where $\sigma_{\text{RJ}}$ is the RMS random jitter. Values range from 0.1~ps (excellent PLL) to 2.0~ps (degraded clock recovery).

\subsection{Block 9: Slicer}

The threshold detector makes a binary decision:
\begin{equation}
    \hat{b}[i] = \begin{cases} 1 & \text{if } r_{\text{sampled}}[i] \geq V_{\text{th}} \\ 0 & \text{otherwise} \end{cases}
\end{equation}
where $V_{\text{th}} = 0$~V for the symmetric NRZ signaling used in UCIe.

%% ====================================================================
\section{Measurement and Analysis (Section 4)}
%% ====================================================================

\subsection{Block 10: Direct BER Measurement}

The bit error rate is computed by direct comparison:
\begin{equation}
    \text{BER} = \frac{1}{N_{\text{bits}}} \sum_{i=1}^{N_{\text{bits}}} \left( b[i] \oplus \hat{b}[i] \right)
\end{equation}

\subsection{Block 11: Voltage Histogram and Q-Factor}

Sampled voltages are separated by the transmitted bit value into two distributions. The statistical Q-factor is:
\begin{equation}
    Q = \frac{\mu_1 - \mu_0}{\sigma_1 + \sigma_0}
\end{equation}
where $(\mu_1, \sigma_1)$ and $(\mu_0, \sigma_0)$ are the mean and standard deviation of the ``1'' and ``0'' voltage distributions, respectively. The estimated BER is then:
\begin{equation}
    \text{BER}_Q = \frac{1}{2} \operatorname{erfc}\!\left(\frac{Q}{\sqrt{2}}\right)
\end{equation}

This provides a statistically robust BER estimate even when the direct error count is zero (i.e., when $N_{\text{bits}}$ is insufficient to observe rare error events)~\cite{bergano1993}.

\subsection{Block 12: Eye Diagram}

The eye diagram is constructed by folding the continuous waveform into 2-UI segments and overlaying up to 1000 traces:
\begin{equation}
    \text{Eye}[k, m] = r_{\text{noisy}}\!\left[(k-1)N + m\right], \quad m = 1, \ldots, 2N
\end{equation}
The resulting overlay provides a visual assessment of eye height (voltage margin), eye width (timing margin), and jitter distribution.

%% ====================================================================
\section{Real-World Test Case Design}
%% ====================================================================

Five test cases are designed to model physically meaningful scenarios encountered in UCIe chiplet integration. Table~\ref{tab:testcases} summarizes their parameters.

\begin{table*}[t]
\caption{Real-World UCIe Link Test Case Parameters}
\label{tab:testcases}
\centering
\begin{tabular}{|l|c|c|c|c|c|c|c|}
\hline
\textbf{Test Case} & \textbf{Data Rate} & $R_{\text{ch}}$ & $C_{\text{ch}}$ & \textbf{FFE} & $\sigma_V$ & $\sigma_{\text{RJ}}$ & \textbf{Physical Scenario} \\
 & (GT/s) & ($\Omega$) & (pF) & ($c_0$, $c_1$) & (mV) & (ps) & \\
\hline
1. Silicon Interposer & 32 & 15 & 0.05 & (1.00, 0.00) & 1 & 0.1 & Ultra-short reach, premium package \\
2. Organic Package & 32 & 40 & 0.15 & (0.85, $-$0.15) & 3 & 0.2 & Standard substrate, moderate loss \\
3. 64 GT/s Future & 64 & 25 & 0.10 & (0.75, $-$0.25) & 5 & 0.25 & Next-gen bandwidth-limited \\
4. Crosstalk Stress & 32 & 30 & 0.10 & (0.90, $-$0.10) & 20 & 0 & Vertical eye closure \\
5. PLL Jitter Failure & 32 & 30 & 0.10 & (0.85, $-$0.15) & 3 & 2.0 & Horizontal eye closure \\
\hline
\end{tabular}
\end{table*}

\subsection{Case 1: Advanced Silicon Interposer}

This case models two chiplets placed micrometers apart on a premium silicon interposer (e.g., TSMC CoWoS or Intel EMIB). The ultra-short trace length results in very low channel parasitics ($R_{\text{ch}} = 15~\Omega$, $C_{\text{ch}} = 0.05$~pF), yielding a channel 3~dB bandwidth of:
\begin{equation}
    f_{3\text{dB}} = \frac{1}{2\pi \times 15 \times 0.05 \times 10^{-12}} \approx 212 \text{ GHz}
\end{equation}
This far exceeds the Nyquist frequency of 16~GHz, resulting in negligible ISI and a wide-open eye requiring no equalization (FFE Preset~0).

\subsection{Case 2: Standard Organic Package}

This case models chiplets connected over a standard organic substrate with longer traces. The increased parasitics ($R_{\text{ch}} = 40~\Omega$, $C_{\text{ch}} = 0.15$~pF) reduce the channel bandwidth to approximately 26.5~GHz, which begins to attenuate the signal's high-frequency content significantly. FFE Preset~2 ($c_0 = 0.85$, $c_1 = -0.15$) provides moderate de-emphasis to counteract this ISI.

\subsection{Case 3: Next-Generation 64 GT/s}

At 64~GT/s, the unit interval shrinks to $T_{\text{UI}} = 15.625$~ps, demanding extreme signal bandwidth. Even with a good-quality channel ($R_{\text{ch}} = 25~\Omega$), the signal experiences severe bandwidth starvation, requiring the strongest FFE Preset~3 ($c_0 = 0.75$, $c_1 = -0.25$) to maintain a usable eye opening.

\subsection{Case 4: Severe Crosstalk and Power Supply Noise}

This stress test models a dense chiplet cluster where adjacent data lanes and noisy power delivery networks inject substantial voltage noise ($\sigma_V = 20$~mV). While the channel itself is benign, the extreme additive noise causes \textbf{vertical eye closure}---the voltage distributions for ``0'' and ``1'' overlap, degrading the Q-factor and increasing BER.

\subsection{Case 5: PLL Clock Jitter Failure}

This case models a degraded Phase-Locked Loop (PLL) with extreme random jitter ($\sigma_{\text{RJ}} = 2.0$~ps), which represents $6.4\%$ of the 31.25~ps UI at 32~GT/s. The voltage noise is kept low ($\sigma_V = 3$~mV) to isolate the timing effect. This jitter causes \textbf{horizontal eye closure}---the receiver samples at incorrect time instants, capturing voltages from the transition regions rather than the center of the eye.

%% ====================================================================
\section{Implementation Details}
%% ====================================================================

\subsection{Software Architecture}

The implementation follows a hierarchical directory structure:

\begin{lstlisting}[caption={Project directory structure},label={lst:dir},language={}]
OELP_project/
  run_ucie_sim.m       % Central runner
  run_test_1.m         % Individual test runners
  run_test_2.m
  run_test_3.m
  run_test_4.m
  run_test_5.m
  run_and_save_test.m  % Auto-save helper
  blocks/
    block_prbs.m
    block_nrz_mapper.m
    block_tx_ffe.m
    block_tx_driver.m
    block_channel.m
    block_rx_load.m
    block_add_noise.m
    block_sample_hold.m
    block_slicer.m
    block_ber_meas.m
    block_histogram_q.m
    block_eye_diagram.m
    plot_block_output.m
    measure_rise_fall.m
    interp_cross.m
\end{lstlisting}

\subsection{Standardized Block Interface}

Every processing block adheres to the interface:
\begin{lstlisting}[caption={Standard block function signature}]
function [output, plotData] = block_xxx(input, cfg)
    if nargin == 0
        % Self-test with default config
        cfg = default_config();
        [input] = upstream_blocks(cfg);
        [output, plotData] = block_xxx(input, cfg);
        plot_block_output(blockNum, plotData);
        return;
    end
    % ... core signal processing ...
end
\end{lstlisting}

The \texttt{plotData} structure returned by each block contains all data needed for visualization, decoupling the signal processing logic from the plotting code.

\subsection{Toolbox Independence}

A key design decision was to eliminate dependency on MATLAB's Signal Processing Toolbox. The channel frequency response computation uses a manually implemented bilinear transform and frequency response evaluation rather than the \texttt{bilinear()} and \texttt{freqz()} functions. This ensures the simulator runs on any standard MATLAB installation.

%% ====================================================================
\section{Results and Discussion}
%% ====================================================================

\subsection{Transmitter Waveform Quality}

The TX driver block produces analog waveforms with measured 20\%--80\% rise times that closely track the target values. At 32~GT/s with a target $t_r = 6.25$~ps, the measured rise time is approximately 7.7~ps after including the output RC filtering effect of the driver impedance. The Gaussian shaping kernel produces smooth, monotonic transitions free of ringing artifacts.

The FFE preset comparison demonstrates clear spectral pre-emphasis: as the de-emphasis ratio increases from Preset~0 to Preset~3, the waveform exhibits progressively sharper transitions at bit boundaries at the cost of reduced signal amplitude during consecutive identical bits.

\subsection{Channel Impact Analysis}

The channel frequency response varies dramatically across test cases. Case~1 (silicon interposer) presents essentially flat response to well beyond the Nyquist frequency, while Case~2 (organic package) exhibits 3--4~dB of attenuation at 16~GHz, causing visible ISI in the time-domain waveform. The channel output plots clearly show the progressive rounding of signal edges as the RC time constant increases.

\subsection{Eye Diagram Observations}

The eye diagrams provide the most intuitive visualization of link quality:

\begin{itemize}
    \item \textbf{Case 1:} Wide-open eye with large voltage and timing margins, confirming that premium silicon interposers require no equalization.
    \item \textbf{Case 2:} Moderately open eye rescued by FFE de-emphasis, demonstrating the equalizer's effectiveness.
    \item \textbf{Case 3:} Narrow eye opening at 64~GT/s despite maximum equalization, illustrating the fundamental bandwidth challenge.
    \item \textbf{Case 4:} Vertically compressed eye due to extreme voltage noise, with overlapping histogram distributions.
    \item \textbf{Case 5:} Horizontally smeared eye due to sampling jitter, with transitions bleeding into the center of the eye.
\end{itemize}

\subsection{Q-Factor and BER Correlation}

The Q-factor provides a reliable predictor of BER across all test cases. For Case~1, the Q-factor exceeds 50, corresponding to negligible BER ($< 10^{-100}$). For the stress cases (4 and 5), the Q-factor degrades to single digits, with the estimated BER rising to observable levels. This validates the statistical framework as a complement to direct error counting.

%% ====================================================================
\section{Conclusion}
%% ====================================================================

This work presents a complete, modular behavioral model for UCIe 3.0 high-speed serial links. The key achievements are:

\begin{enumerate}
    \item A 12-block architecture that maps one-to-one with the physical signal flow from transmitter silicon through the channel to the receiver decision circuit.
    
    \item Parameterized channel modeling that captures the distinction between premium silicon interposers and standard organic substrates through configurable RC values.
    
    \item Separate modeling of vertical eye closure (voltage noise) and horizontal eye closure (timing jitter), enabling independent analysis of noise and jitter budgets.
    
    \item A fully automated test infrastructure with per-block plotting and PNG export for documentation and reporting.
    
    \item Toolbox-independent implementation using only core MATLAB functions, ensuring broad accessibility.
\end{enumerate}

The modular architecture readily extends to more advanced equalization schemes (e.g., CTLE, DFE), multi-level signaling (PAM-4), and S-parameter-based channel models. Future work will incorporate measured channel data from physical UCIe test vehicles and validate the behavioral model against silicon measurements.

\begin{thebibliography}{00}

\bibitem{ucie_spec} Universal Chiplet Interconnect Express (UCIe) Specification, Rev.~1.1, UCIe Consortium, 2023.

\bibitem{razavi2012} B.~Razavi, ``The StrongARM latch,'' \textit{IEEE Solid-State Circuits Magazine}, vol.~7, no.~2, pp.~12--17, 2015.

\bibitem{fan2007signal} J.~Fan, X.~Ye, J.~Kim, B.~Archambeault, and A.~Orlandi, ``Signal integrity design for high-speed digital circuits: Progress and directions,'' \textit{IEEE Trans. Electromagn. Compat.}, vol.~52, no.~2, pp.~392--400, 2010.

\bibitem{jeruchim2000} M.~C.~Jeruchim, P.~Balaban, and K.~S.~Shanmugan, \textit{Simulation of Communication Systems}, 2nd~ed. New York: Kluwer, 2000.

\bibitem{proakis2008} J.~G.~Proakis and M.~Salehi, \textit{Digital Communications}, 5th~ed. New York: McGraw-Hill, 2008.

\bibitem{bergano1993} N.~S.~Bergano, F.~W.~Kerfoot, and C.~R.~Davidson, ``Margin measurements in optical amplifier systems,'' \textit{IEEE Photon. Technol. Lett.}, vol.~5, no.~3, pp.~304--306, 1993.

\end{thebibliography}

\end{document}
